\section{Hạn chế và hướng phát triển}
\label{sec:limitations}

Báo cáo có 11 hạn chế. Hạn chế~\ref{lim:etth2} và \ref{lim:pipeline} rút ra khi phân tích ETTh2 và số công bố; các hạn chế còn lại đã biết từ khi thiết kế thí nghiệm.

\begin{enumerate}[label=\arabic*., ref=\arabic*]
    \item \label{lim:data}\textbf{Dữ liệu.} Chỉ ba dataset ETT, cùng 7 kênh, và chỉ là hai nguồn độc lập: ETTh1 và ETTm1 đo cùng một máy biến áp, chỉ khác tần suất lấy mẫu; ETTh2 đo máy thứ hai. Một kết quả lặp lại trên cả ETTh1 lẫn ETTm1 vì vậy không phải hai lần xác nhận độc lập. Phép kiểm qua module LTSF-Linear và phép quét kernel của RQ1 chỉ chạy ở ETTh1, $H = 96$. Dẫn xuất đại số đúng cho mọi ô, nhưng độ lớn của phần tách theo $k$ chưa đo ở ô khác. \emph{Hướng phát triển:} thêm các dataset khác của LTSF-Linear (Electricity, Traffic, Weather) và quét kernel trên toàn ma trận.
    \item \label{lim:models}\textbf{Mô hình.} Chỉ mô hình tuyến tính, dùng chung trọng số cho mọi kênh, với $L = 336$ cố định và $k = 25$ trong grid. Kết luận không áp dụng trực tiếp cho biến thể \texttt{individual = True} hay cho mô hình phi tuyến.
    \item \label{lim:valtest}\textbf{Val và test của ETT thuộc hai chế độ khác nhau.} Trên ETTh1, MSE val cao hơn hẳn MSE test. Mọi lựa chọn dựa trên val ($\lambda^*$, lr, epoch dừng) chịu giới hạn này: quy tắc ba trường hợp gán nhãn ``khác biệt do val khác test'' cho \nmBFourValTestCount/36 ô với lr gốc và \nmTunedValTestCount/36 ô với lr tinh chỉnh (\tabref{tab:b4}). Ví dụ khác: trên ETTh1 và ETTm1, co ridge về dự báo lặp giá trị cuối thay vì về $0$ làm MSE val thấp hơn nhưng MSE test cao hơn ở gần như mọi tổ hợp đã thử. \emph{Hướng phát triển:} chọn siêu tham số bằng kiểm định chéo theo thời gian (nhiều đoạn val) thay vì một đoạn val duy nhất.
    \item \label{lim:outlier}\textbf{Outlier nhân tạo.} Chỉ một dạng outlier: cộng $\pm k$ độ lệch chuẩn vào từng điểm ở vị trí ngẫu nhiên, chỉ trên ETTh1, chỉ với Linear và NLinear. Thí nghiệm độ lớn ($3\sigma$, $5\sigma$) chỉ có Linear, $H \in \{96, 720\}$ và một seed. Outlier thật (như các đoạn cảm biến bị kẹt của ETTh2) có cấu trúc khác.
    \item \label{lim:delta}\textbf{Một $\delta$ chung cho cả ô.} $\hat\sigma$ của phần dư theo từng bước $h$ chênh nhau nhiều trong một ô, nên một $\delta$ chung nằm trong vùng bậc hai với các bước gần và gần vùng tuyến tính với các bước xa. Chuẩn hóa phần dư theo $\hat\sigma_h$ rồi dùng một $\delta$ chung (tức $\delta_h = 1{,}345\,\hat\sigma_h$) chỉ được thử ở \nmDeltaHCellCount{} ô (ETTh2 $H = 720$ và ETTh1 $H = 96$, Linear). Ở cả hai ô, Huber vẫn đứng giữa MSE và MAE (\tabref{tab:deltah}), nhưng kết luận này chưa được kiểm trên toàn ma trận.
    \item \label{lim:lr}\textbf{Lưới learning rate.} Lưới lr ban đầu là $\times 0{,}5$ đến $\times 8$ lr gốc; \nmTunedLrEdgeOldCount/\nmTunedLrCellCount{} ô chọn mức ở đầu mút. Ở các ô này, lưới được mở rộng thêm tối đa hai mức theo cùng một quy tắc ($\times 0{,}25$, $\times 0{,}125$ hoặc $\times 16$, $\times 32$), và lr chọn đổi ở \nmTunedLrExtChangedCount{} ô. \nmTunedLrEdgeCount{} ô vẫn chọn mức xa nhất đã thử (\nmTunedLrStillEdgeList), nên lr tối ưu trên val của chúng có thể còn nằm ngoài khoảng đã thử. Lựa chọn lr cũng nhiễu (\nmTunedLrNoisyCount/\nmTunedLrCellCount{} ô có khoảng cách tới mức thứ hai dưới một sai số chuẩn). Chỉ lr được tinh chỉnh; batch, số epoch, patience giữ như script gốc.
    \item \label{lim:etth2}\textbf{Cơ chế trên ETTh2 chưa rõ.} SGD tốt hơn nghiệm tối ưu trên cả val và test ở ETTh2 với Linear và DLinear. Các giải thích đã kiểm (co về trọng số khác $0$, đoạn cảm biến bị kẹt, dừng sớm của gradient descent toàn batch) đều không khớp. Giả thuyết nhiễu đầu vào ở \secref{sec:discussion} cũng chưa được kiểm.
    \item \label{lim:pipeline}\textbf{Pipeline SGD và số công bố.} SGD của repo khác code gốc của LTSF-Linear ở vài chi tiết (cách tính val, thứ tự lấy số ngẫu nhiên). Code gốc đã được chạy ở ETTh2 $H \in \{336, 720\}$ (Linear, DLinear), ở ETTm1 $H = 96$ DLinear trên CPU (MSE test \nmOffETTmOneHNineSixMean{} $\pm$ \nmOffETTmOneHNineSixSd, \nmOffETTmOneHNineSixN{} seed, so với số công bố \nmBOneETTmOneHNineSixDLinearPub), và ở cả \nmDropCellCount{} ô có số công bố thấp hơn mọi seed của repo, với \texttt{drop\_last} của test loader bật và tắt (\tabref{tab:droplast}). Môi trường của chúng tôi (PyTorch 2.14.0, CUDA 13.0, GPU RTX 5060 Ti hoặc CPU, Windows 11) khác môi trường của bài gốc, và mỗi ô chỉ có \nmDropSeedCount{} seed với code gốc. Với \texttt{drop\_last=True}, số công bố nằm trong phân phối theo seed của code gốc ở \nmDropPubInTrueCount/\nmDropCellCount{} ô; ở các ô còn lại chúng tôi không đạt được số công bố, và \texttt{drop\_last} không giải thích được khoảng chênh. Chúng tôi không kết luận số công bố là không tái lập được: khoảng chênh có thể đến từ môi trường, từ số seed hạn chế, hoặc từ chi tiết khác của lần chạy gốc mà repo không ghi lại.
    \item \label{lim:irls}\textbf{Độ chính xác của IRLS.} IRLS dừng theo ngưỡng giảm tương đối của hàm mục tiêu, nên nghiệm MAE là xấp xỉ. Trên ETTh1 $H = 96$, nghiệm từ ba điểm khởi tạo có cùng hàm mục tiêu tới sai số rất nhỏ, nhưng trọng số còn khác nhau do dừng sớm (Phụ lục~\ref{app:irls}). MAE được tính qua Huber với $\delta = 10^{-6}$, tức một bản làm trơn của MAE.
    \item \label{lim:metrics}\textbf{Metric trên dữ liệu đã chuẩn hóa.} Như LTSF-Linear, mọi MSE và MAE tính trên dữ liệu đã chuẩn hóa theo train (\secref{sec:method-data}), nên không đọc trực tiếp được theo đơn vị gốc của từng kênh, và các kênh có độ lệch chuẩn khác nhau được tính với trọng số như nhau.
    \item \label{lim:commits}\textbf{Commit ghi trong kết quả.} Một số JSON ghi commit có hậu tố \texttt{-dirty}, vì lúc chạy có file script hoặc báo cáo chưa commit. Code của solver và mô hình không đổi ở các commit đó; danh sách file cụ thể ghi trong từng báo cáo kỹ thuật ở \texttt{docs/reports/} (Phụ lục~\ref{app:repro}).
\end{enumerate}
